\documentclass[hyperref={pdfpagelabels=false},table, handout]{beamer}

\input{lec_common}
\begin{document}
\part{Unit 2:Strings, Booleans, Functions and Plot.}


\begin{frame}[fragile]{In this lecture}

\begin{itemize}
  \item What is Scientific Computing
  \item More about lists and operators
  \item Calculating sums
  \item Conditional statements
  \item More about for loops
  \item String formatting
  \item Functions
  \item Modules
  \item More about plots
  \item Fixed point iterations
\end{itemize}
\end{frame}

%%%%%%%%%%%%%%%%%%%%%%%%%%%%%%%%%%%%%%%%%%%%%%%%%%%%%%%%%%%%%%%%%%%%%
%%%%%%%%%%%%%%%%%%%%%%%%%%%%%%%%%%%%%%%%%%%%%%%%%%%%%%%%%%%%%%%%%%%%%
\begin{frame}\frametitle{Scientific Computing 1/2}

  \small

  \structure{Computational science}, also known as \structure{scientific computing} or \structure{scientific computation} (SC),
  is a rapidly growing field that uses advanced computing capabilities to understand and solve complex problems.
  It is an area of science which spans many disciplines, but at its core,
  it involves the development of models and simulations to understand natural systems (\url{en.wikipedia.org/wiki/Computational_science}). \\

  \vskip 5pt

  \begin{itemize}
    \item \structure{Algorithms} (numerical and non-numerical): mathematical models, computational models, and computer simulations developed to solve science (e.g., biological, physical, and social), engineering, and humanities problems
    \item \structure{Computer hardware}: that develops and optimizes the advanced system hardware, firmware, networking, and data management components needed to solve computationally demanding problems
    \item \structure{The computing infrastructure} that supports both the science and engineering problem solving and the developmental computer and information science.  \\
  \end{itemize}

  \alert{Note}: \\
  Scientific Computing (Computational Science) $\neq$ Computer Science
\end{frame}

%%%%%%%%%%%%%%%%%%%%%%%%%%%%%%%%%%%%%%%%%%%%%%%%%%%%%%%%%%%%%%%%%%%%%
%%%%%%%%%%%%%%%%%%%%%%%%%%%%%%%%%%%%%%%%%%%%%%%%%%%%%%%%%%%%%%%%%%%%%
\begin{frame}\frametitle{Scientific Computing 2/2}

Because we want to solve large (many unknowns) problems with computer simulations
we need fast algorithms but also fast and efficient software.
This is the main difference to other areas of programming. \\

\vskip 10pt

We often have to make sacrifices in design and code style in order to achieve performance.
This leads to programs with millions of lines of code that are difficult to understand and manage.

\vskip 10pt
  Example: The Community Earth System Model (\url{https://www.cesm.ucar.edu})

\end{frame}

%%%%%%%%%%%%%%%%%%%%%%%%%%%%%%%%%%%%%%%%%%%%%%%%%%%%%%%%%%%%%%%%%%%%%
%%%%%%%%%%%%%%%%%%%%%%%%%%%%%%%%%%%%%%%%%%%%%%%%%%%%%%%%%%%%%%%%%%%%%
\begin{frame}[fragile]\frametitle{Revision -- lists}

  Given a list \pyth{L} you can use the function \pyth{len(L)} to find the length of the list.

\begin{python}
  L = [4, 9, -1, 2]
  print("The length is", len(L)) # The length is 4
\end{python}

  \pause

  \vskip 5pt

  You can also use the method \pyth{append} by using dot notation after the name of the list.

\begin{python}
  L.append( 11 )
  print("L = ", L) # L = [4, 9, -1, 2, 11]
\end{python}


\end{frame}

%%%%%%%%%%%%%%%%%%%%%%%%%%%%%%%%%%%%%%%%%%%%%%%%%%%%%%%%%%%%%%%%%%%%%
%%%%%%%%%%%%%%%%%%%%%%%%%%%%%%%%%%%%%%%%%%%%%%%%%%%%%%%%%%%%%%%%%%%%%
\begin{frame}[fragile]\frametitle{Some more functions}

 Given a list \pyth{L} you can also use:

  \begin{itemize}
    \item \pyth{sum( L )}
    \item \pyth{max( L )}
    \item \pyth{min( L )}
  \end{itemize}

  \pause

\begin{python}
  L = [4, 9, -1, 2]
  print("The sum is", sum(L)) # The sum is 14
  print("The max is", max(L)) # The max is 9
  print("The min is", min(L)) # The min is -1
\end{python}

\end{frame}

%%%%%%%%%%%%%%%%%%%%%%%%%%%%%%%%%%%%%%%%%%%%%%%%%%%%%%%%%%%%%%%%%%%%%
%%%%%%%%%%%%%%%%%%%%%%%%%%%%%%%%%%%%%%%%%%%%%%%%%%%%%%%%%%%%%%%%%%%%%
\begin{frame}[fragile]\frametitle{Even more functions}

 Given a list \pyth{L} you can also use:

  \begin{itemize}
    \item \pyth{L.reverse()}
    \item \pyth{L.sort()}
    \item \pyth{L.sort(reverse=True|False, key=myFunc)}
  \end{itemize}

  \pause

\begin{python}
  L = [4, 9, -1, 2]
  print("L = ", L)  # L = [4, 9, -1, 2]
  L.reverse()
  print("L = ", L)  # L = [2, -1, 9, 4]
  L.sort()
  print("L = ", L)  # L = [-1, 2, 4, 9]
\end{python}

\end{frame}

%%%%%%%%%%%%%%%%%%%%%%%%%%%%%%%%%%%%%%%%%%%%%%%%%%%%%%%%%%%%%%%%%%%%%
%%%%%%%%%%%%%%%%%%%%%%%%%%%%%%%%%%%%%%%%%%%%%%%%%%%%%%%%%%%%%%%%%%%%%

\begin{frame}[fragile]{Programming example}
\vfill
How do we calculate
\[
\sum_{i=0}^n i \ ?
\]

\pause

  Discussion (5min).

\vskip 10pt

\pause

\begin{python}
  # for loop
  s = 0
  for i in range(0,101): # n = 100
      s += i
  print( s ) # 5050
\end{python}

\pause

\begin{python}
  # list comprehension
  L = [i for i in range(0,101)]
  print( sum(L) ) # 5050
\end{python}
\end{frame}

%%%%%%%%%%%%%%%%%%%%%%%%%%%%%%%%%%%%%%%%%%%%%%%%%%%%%%%%%%%%%%%%%%%%%
%%%%%%%%%%%%%%%%%%%%%%%%%%%%%%%%%%%%%%%%%%%%%%%%%%%%%%%%%%%%%%%%%%%%%

\begin{frame}[fragile]{Conditional Expressions}
\book{20, 22, 37}
\begin{defblock}A \alert{conditional expression} is an expression that may have the \alert{Boolean} value \alert{True} or \alert{False}.
\end{defblock}
\pause
Some common operators that yield conditional expressions are:
\begin{itemize}
\item \pyth{==}, \pyth{!=}
\item \pyth{<}, \pyth{>}, \pyth{<=}, \pyth{>=}
\item One combines different Boolean values with \pyth{or} and \pyth{and}
\item \pyth{not} gives the \alert{logical negation} of the expression that follows
\end{itemize}
\end{frame}

\begin{frame}[fragile]{Conditional Expression Examples}
\begin{example}
\begin{python}
2 >= 4  # False
2 < 3 < 4 # True
2 < 3 and 3 < 2 # False
2 != 3 < 4 or False # True
2 <= 2 and 2 >= 2 # True
not 2 == 3 # True
not False or True and False # True!
\end{python}
\end{example}
Note in the last example the priority rules when using  \pyth!not!, \pyth!and!, and \pyth!or!.\\
They can be compared to addition ($\rightarrow$ \pyth!or!), multiplication, ($\rightarrow$~\pyth!and!),
and sign change ($\rightarrow$ \pyth!not!) in classical arithmetics.


\end{frame}

\begin{frame}[fragile]{Conditional statements}
\begin{alertblock}{\pyth!If! statement}
A conditional statement delimits a \alert{block} that will be executed if the condition is true. An \alert{optional} block, started with the keyword \pyth{else} will be executed if the condition is not fulfilled.
\end{alertblock}
\begin{example}%{Example: Absolute Value}
We print the absolute value of $x$:\\[0.5cm]
\begin{minipage}[c]{0.4\textwidth}
Mathematics:\\
\[
|x| = \left \{
\begin{aligned}
 x & \text{  if } &x \ge 0 \\
-x & \text{  else } &
\end{aligned}
\right .
\]
\vfill
\end{minipage}
\hfill $\Longrightarrow$ \hfill
\begin{minipage}[c]{0.4\textwidth}
Python:
\begin{python}
x = ...
if x >= 0:
  print(x)
else:
  print(-x)
\end{python}
\end{minipage}
\end{example}
\end{frame}

\begin{frame}[fragile]{Error messages}
By using Boolean expressions you can check for potential errors:
\begin{python}
denom = ...
num = 3
if not denom == 0:
   frac = num / denom
else:
   raise ZeroDivisionError("Don't divide by zero!")
\end{python}

The string in the parenthesis is up to the programmer. \\
The error type has to be predefined. \\ \hfill {\footnotesize See later chapter}.
\end{frame}




\begin{frame}[fragile]{The full \pyth!for! statement: \texttt{break}}
\book{21, 22}

%The for statement has two important keywords: \pyth{break} and \pyth{else}.
%\begin{itemize}[<+->]
%\item
\pyth!break! gets out of the for loop even if the list we are iterating is not exhausted.
\begin{python}
for x in x_values:
  if x > threshold:
    break
  print(x)
\end{python}
%\item
\end{frame}
\begin{frame}[fragile]{The full for statement: \pyth!else!}
 \alert{else} checks whether the for loop was \emph{broken} with the \pyth!break! keyword.
\begin{python}
for x in x_values:
  if x > threshold:
    break
else:
  print("all the x are below the threshold")
\end{python}
If we \emph{did not break} the \pyth{else}-block is executed.
%\end{itemize}
\end{frame}



\begin{frame}[fragile]{Basic string formatting}
\book{43}
\begin{itemize}
\item Using keywords:
\begin{python}
"I'm a {something}.".format(something="string")
# "I'm a string."
"I'm a {something}.".format(something=10)
# "I'm a 10."
\end{python}

\item Using positional arguments:
\begin{python}
"Two strings, {} and {}.".format('a', 'b')
# 'Two strings, a and b.'
"Pi = {0} and e = {1}.".format(pi, e)
# 'Pi = 3.14159265359 and e = 2.71828182846.'
\end{python}

\item Numbers can be formatted:
\begin{python}
quantity = 33.45
"{0:f}".format(quantity)    # 33.450000
"{0:1.1f}".format(quantity) # 33.5
"{0:.2e}".format(quantity)  # 3.35e+01
\end{python}
\end{itemize}

\end{frame}
\begin{frame}[fragile]{Basic string formatting \hfill (Cont.)}
A way to work with string formatting is to construct a string
as a template and then applying the format method in a second step:
\vfill
\begin{python}
myname="I'm {name}. I work as a {profession}."
print(myname.format(name='Erika', profession='manager')
print(myname.format(name='Carl-Gustaf', profession='king')
\end{python}
\vfill

\end{frame}
\begin{frame}[fragile]{f-strings}
An easy and direct alternative is the use of f-strings:
\vfill
\begin{python}
result = 2.54
p=(f'The pressure of the tire is {result} bar\n'
   f'which is {result-2.5:1.2f} bar too much!')
print(p)
\end{python}
\vfill
This results in:
\begin{verbatim}
The pressure of the tire is 2.54 bar
which is 0.04 bar too much!
\end{verbatim}
{\tiny Note, how calculations and formatting can be done within an f-string.\\ Note also the implicit continuation line forced by the use of the parantheses.}
\end{frame}
\begin{frame}[fragile]{Old string formatting (just for reference)}
Older versions of Python use the following syntax to format strings. It works in newer versions too.
\begin{itemize}
\item for strings:
\begin{python}
course_code = "NUMA21"
print("This course's name is %s" % course_code)
# This course's name is NUMA21
\end{python}
\item for integers:
\begin{python}
nb_students = 16
print("There are %d students" % nb_students)
# There are 16 students
\end{python}
\item for reals:
\begin{python}
average_grade = 3.4
print("Average grade: %f" % average_grade)
# Average grade: 3.400000
\end{python}
\end{itemize}
\end{frame}



\begin{frame}[fragile]{Basics on Functions}
\book{23}
Consider the following mathematical \alert{function}:
\[x \mapsto f(x) := 2x + 1\]
The Python equivalent is:
\begin{python}
def f(x):
  return 2*x + 1
\end{python}
\begin{itemize}
\item the keyword \pyth{def} tells Python we are defining a function
\item \texttt{f} is the \alert{name} of the function
\item \texttt{x} is the \alert{parameter}, or \alert{input} of the function
\item what follows \pyth!return! is called the \alert{output} of the function
\end{itemize}
\end{frame}


\begin{frame}[fragile]{Calling a Function}
Once the following function is defined:
\begin{python}
def f(x):
  return 2*x + 1
\end{python}
it may now be called using:
\begin{python}
f(2) # 5
f(1) # 3
# etc.
\end{python}

Note, the \alert{parameter} $x$ is replaced by the \alert{argument} $2$.
\end{frame}



\begin{frame}[fragile]{Writing scripts (Book p. 15)}

Collect a sequence commands in a file with its name ending in \alert{.py}, e.g.\ \pyth{smartscript.py}
\begin{python}
def f(x):
  return 2*x + 1

z = []
for x in range(10):
  if f(x) > pi:
    z.append(x)
  else:
    z.append(-1)
print(z)
\end{python}

\pause

From IPython:
\begin{python}
run smartscript.py # [-1, -1, 2, 3, 4, 5, 6, 7, 8, 9]
\end{python}

\pause

From command prompt (i.e. in the terminal window)
\begin{python}
python smartscript.py
\end{python}

\end{frame}



\begin{frame}[fragile]{Example - collecting functions}
\book{24}
Create a \alert{module} by collecting functions into a single file, e.g.\ \pyth{smartfunctions.py} as:
\begin{python}
def g(x):
  return x**2 + 4 * x - 5

def h(x):
  return 1/f(x)

def f(x):
  return 2 * x + 1
\end{python}
\begin{itemize}
\item These functions can now be used by any external script or directly in the IPython environment.
\item Functions within the module can depend on each other.
\item Grouping functions with a common theme or purpose gives modules that can be shared and used by others.
\end{itemize}
\end{frame}


\begin{frame}[fragile]{Using modules \hspace*{\fill}}
\book{25}
Modules need to be imported.
\begin{python}
import smartfunctions
print(smartfunctions.f(2)) # 5

from smartfunctions import g  #import just only g
print(g(1)) # 0

from smartfunctions import * #import all
print(h(2) * f(2)) # 1.0
\end{python}
{\small Note the use of the namespace \pyth{smartfunctions}!}
\end{frame}




\begin{frame}[fragile]{Plot}
\book{123}
\begin{itemize}
\item
You have already used the command \pyth{plot}. It needs a list of $x$ values and a list of $y$ values. If a single list, $y$ is given, the list \pyth{list(range(len(y)))} is assumed as $x$ values.
\item You may use the keyword argument \pyth{label} to give your curves a name, and then show them using \pyth{legend}.
\begin{python}
x_vals = [.2*n for n in range(20)]
y1 = [sin(.3*x) for x in x_vals]
y2 = [sin(2*x) for x in x_vals]
plot(x_vals ,y1, label='0.3')
plot(x_vals, y2, label='2')
legend()
show() # not always needed.
\end{python}
\end{itemize}
\end{frame}


\begin{frame}[fragile]{Plot (Cont.)}
Here comes the plot with a legend.
\begin{center}
 \includegraphics[height=3cm]{legend_ex.pdf}
 % legend_ex.pdf: 586x442 pixel, 72dpi, 20.67x15.59 cm, bb=0 0 586 442
\end{center}
Here is an example with more keywords:
\begin{python}
 plot(x_vals,y2,
      color='green',
      linestyle='dashed',
      marker='o',
      markerfacecolor='blue',
      markersize=12,linewidth=6)
\end{python}
\end{frame}
\begin{frame}[fragile]{Plot (Cont.)}
... and here the resulting graph:
\vfill
\begin{center}
 \includegraphics[height=6cm]{legend_ex2.pdf}
 % legend_ex.pdf: 586x442 pixel, 72dpi, 20.67x15.59 cm, bb=0 0 586 442
\end{center}
\vfill
\end{frame}


\end{document}
